\section{引言}

海洋大数据的计算引擎正从岸基超算中心向移动观测平台边缘延伸，使“继续计算、补充观测还是立即行动”成为自主水下航行器必须面对的在线决策问题。传统流程依赖岸基中心汇集卫星与模式产品并下发指令\cite{chassignet2007hycom,chin2017mur}，但水下声学通信带宽极窄且上浮卫星通信耗时耗能，长航程航行器必须依靠低功耗机载计算模块独立完成环境推演与重规划\cite{remus100,jetsonorinnx}。由于机载算力与电池能量受限，航行器既不能无限制地进行高分辨率搜索，也不能频繁中断航行去原位测流或上浮通信\cite{horvitz1987reasoning,russell1991right}。如何在计算、观测与行动之间分配资源，构成了海洋边缘智能的核心挑战\cite{sung2021quit,magnino2026metareasoning}。

在台湾以东至东海的黑潮强流区，当航行器携带数天前下发的粗分辨率海流图执行绕障航行时，选错绕行方向将导致数小时的航时损失。长航程航行器对水巡航速度约为1.0米/秒，与黑潮主轴及邻近涡旋流速相当。在相距百余千米的起终点之间绕越禁航区时，从左侧或右侧绕行是不可逆的二元拓扑决策：一旦驶入逆流一侧，不仅航时大增，还可能因顶流过强而无法通行。本文考虑航行器只携带3天前制作的1°粗分辨率海流图（简称粗图）的情形：该图既丢失了中尺度锋面与涡旋切变，又包含数天内流轴摆动的时效误差，极易误导左右选边。

面对粗图上的航路决策，常规做法是在当前海流图上投入更多算力以加密网格与航向离散度\cite{lolla2014levelset}，但当环境信息本身缺失时，这种确定性等价做法难以奏效。若左右两侧在粗网格上的预计航时接近，规划器会自然地细分网格以分辨快慢。然而，这混淆了求解器粗糙与海洋环境不确定性：加密网格只能减少同一张粗图上的离散误差，却不能提供任何关于真实涡旋位置的新信息\cite{jiang2015horizon,cuadron2025overthinking}。正如第3节前期赤潮巡测研究所示，当背景场陈旧时，成百上千倍地增加机载推演只会使规划器在想象模型上的自评不断虚高，而航线在真实海洋中的价值早已饱和\cite{smith2006curse}。

第二种常见路径是利用集合预报开展风险感知航路规划，但现有方法不仅无法回答何时应停止计算或去观测，还容易因有限抽样漏掉翻转决策的流场\cite{wang2016path,subramani2018stochastic,kularatne2018forecast}。集合规划通过在数十个扰动流场中评估平均航时或翻转比例来折中选路，但将不确定性视为静态输入，缺乏在计算与观测之间切换的元推理机制\cite{lin2015metareasoning,elboher2023formal}。进一步而言，集合给出的翻转比例只反映有多少成员翻转了决定，而不反映翻转后损失有多大，因而难以区分“翻转后多花半小时”与“翻转后多花十几个小时”的任务，容易低估选错一侧的尾部风险\cite{benhaim2006infogap,marx2020predictive}。

第三种路径是基于信息价值的自适应采样，即主动前往不确定性最高的海区测量，但对整个海流场有价值的信息并不等同于对左右选边有价值的信息\cite{fossum2019,hollinger2014sampling}。传统自适应采样以降低全场估计方差或最大化互信息为目标\cite{srinivas2012gpucb}；然而在绕障任务中，若禁航区外某处的误差对左右航时影响相同，测量它便徒耗航时；反之，顺路在起点附近测流虽代价低廉，若决定两侧快慢的关键切变远在数十千米之外，近处测量也无法改变决策。序贯实验设计与模型判别理论指出，真正决定是否需要新信息的，是在现有证据下能否找到与当前观测同样相容、却要求采取相反行动的竞争模型\cite{chernoff1959sequential,box1967discrimination,garivier2016optimal,attia2018goal}。

为此，本文提出“合理反例检验”：机载计算与现场观测消除的是两类不同性质的不确定性；当两侧难分伯仲是因为求解网格粗糙（数值歧义）时，多算即可做对决定；而当存在与粗图同样相容、却使另一侧明显更快的合理海流（环境歧义）时，继续计算无法分辨真伪，唯有引入新观测才能消解，且在观测过贵时应直接选择最坏后悔更小的一侧。我们将这一判别转化为操作问题：针对粗图推荐的一侧，能否在多年海洋大数据标定的合理集合内搜出使另一侧更快、且让当前选择付出沉重代价的合理反例？在图\ref{fig:overview}所示的2022年1月18日测试任务示例中，粗图上两侧只差约0.3小时且地图推荐左侧，继续加密网格只会把这一账面差值算得更细，而合理反例检验在合理集合内搜出了使左侧明显变慢、右侧更快的物理相容流场，预示了真实海洋中右侧快11.7小时的结果，并据此触发门控决策。

\begin{figure}[htbp]
\centering
\includegraphics[width=\linewidth]{fig_route_overview}
\caption{粗分辨率海流图下的航路决策与合理反例检验（2022年1月18日任务示例）。(a) 航行器携带的3天前1°粗图，左侧航时37.7小时、右侧38.0小时，地图选左侧；(b) 在大数据标定的合理集合内搜出的合理反例海流，左侧55.7小时、右侧44.1小时；(c) 当天1/12°真实海流，左侧53.4小时、右侧41.7小时，右侧实际快11.7小时；(d) 本文决策流程：先校验粗图两侧时间差是否小于数值误差容限以决定是否按需精算，再通过反例检验计算最坏后悔，最坏后悔小则照地图走，大则在观测便宜时上浮取新图、观测昂贵时改选最坏后悔更小的保守一侧。再多的机载计算也只能把粗图上的0.3小时差值算准，唯有合理反例检验能显式找出(b)(c)这类同样合理却翻转决策的海洋状态。}
\label{fig:overview}
\end{figure}

本文的主要贡献如下：
\begin{itemize}
    \item \textbf{概念与判别准则创新：}据我们所知，本文首次在自主水下航行器航路决策中分离了可由计算消除的“数值歧义”与需靠新信息或保守决策应对的“环境歧义”，提出以最坏合理后悔为核心判据的合理反例检验框架，明确了边缘计算与海洋观测的功能边界。
    \item \textbf{大数据标定与高效反例搜索：}利用多年再分析海流标定旧粗图的无散流函数误差模态及覆盖90\%以上真实时间差的合理集合半径，结合包络定理梯度与球面投影搜索，以约71次标准网格求解直接定位最坏后悔流场，其选边效果与按期望时间准则使用约1024个随机合理场（2048次求解，约为其29倍）相当。
    \item \textbf{黑潮、湾流与阿古拉斯回流区三个强流区的系统实证：}在六年黑潮再分析海流共1333个任务上的实验表明，64倍算力仅使选错率从21.7\%降至19.7\%，而99.7\%的任务存在合理反例；最坏后悔前20\%的任务集中了67\%的损失，门控后取保守将平均航时损失降低28\%。在湾流区以及在下载数据前即写定检验计划的阿古拉斯回流区，门控识别高危任务与门控后上浮取图（平均损失分别从3.56小时降至1.31小时、从6.87小时降至2.20小时）同样成立，而不观测的保守改选及其改进规则均无显著收益，从而表明门控的核心价值在于决定何时去观测；黑潮区改带真实卫星测高图时，加算无效与风险排序仍成立，而各种响应的收益不再显著。
\end{itemize}

本文其余部分组织如下：第2节回顾相关工作；第3节简述前期赤潮巡测研究；第4节建立绕障航路决策模型；第5节详述合理反例检验与门控方法；第6节给出实验结果、与随机合理场的比较、适用范围与两个海区的重复实验；第7节拆解创新点与研究脉络；第8节总结全文。
